The first line contains a single integer $t$ ($1 \le t \le 3 \cdot 10^5$) --- the number of test cases.

The first line of each test case contains two integers $n$ and $q$ ($2 \le n \le 3 \cdot 10^5$, $1 \le q \le 3 \cdot 10^5$) --- the number of vertices on the ring and the number of events.

The second line contains $n$ integers $a_1, a_2, \ldots, a_n$ ($-10^9 \le a_i \le 10^9$) --- the weights of the ring edges: edge $i$ connects vertices $i$ and $(i \bmod n) + 1$.

The third line contains $n$ integers $b_1, b_2, \ldots, b_n$ ($-10^9 \le b_i \le 10^9$) --- the weights of the edges between vertex $i$ and the central vertex.

Each of the next $q$ lines describes one event in one of the following formats:
\begin{itemize}
\item \texttt{1 s t} ($1 \le s < t \le n+1$, where vertex $n+1$ is the central vertex) --- print the cheapest valid route from vertex $s$ to vertex $t$;
\item \texttt{2 x w} ($1 \le x \le n$, $-10^9 \le w \le 10^9$) --- set $a_x := w$;
\item \texttt{3 x w} ($1 \le x \le n$, $-10^9 \le w \le 10^9$) --- set $b_x := w$.
\end{itemize}

It is guaranteed that the sum of $n$ over all test cases does not exceed $3 \cdot 10^5$, and that the sum of $q$ over all test cases does not exceed $3 \cdot 10^5$.
